\documentclass[10pt,a4paper]{amsart}
\usepackage[margin=19mm]{geometry}
\usepackage{amsmath,amssymb,mathtools}
\usepackage[scr=rsfs,cal=euler]{mathalfa}
\usepackage{xfrac}
\usepackage[unicode,hidelinks,bookmarks=false]{hyperref}
\newcommand{\ie}{i.e.\ }
\newcommand{\cf}{cf.\ }
\newcommand{\resp}{respectively\ }
\newcommand{\Subsection}{\subsection}
\providecommand{\nobreakdash}{\nobreak\hbox{-}}
\setlength{\emergencystretch}{2em}
\multlinegap0pt
\usepackage{mathtools}
\usepackage{amssymb}
\usepackage{float}
\usepackage{xcolor}
\usepackage{algorithm}\usepackage{algpseudocode}
\usepackage{tikz}
\usepackage{bm}
\usepackage{dsfont}
\usepackage{csquotes}
\newtheorem{theorem}{Theorem}
\title{\textbf{Block-Separated Overpartitions: Fibonacci Structure and Euler Factorization}}
\author{El-Mehdi Mehiri}
\date{Source: arXiv:2511.16580v2 (CC BY 4.0)}
\begin{document}
\maketitle

\noindent\emph{Source and licence.} \textbf{Block-Separated Overpartitions: Fibonacci Structure and Euler Factorization}. Version arXiv:2511.16580v2 (CC BY 4.0), distributed by arXiv under the Creative Commons Attribution 4.0 International licence, http://creativecommons.org/licenses/by/4.0/, as recorded in the arXiv licence metadata for this exact version and shown on the abstract page https://arxiv.org/abs/2511.16580v2. Full licence notice: \texttt{licenses/arxiv-2511.16580-CC-BY-4.0.txt}.\par\smallskip

\noindent\textit{Editorial scope.} Counted unit: the article's theorem thm:euler-factorization (source0.tex lines 762--776). The article's complete original proof of it is delivered with it (source0.tex lines 777--876, one \texttt{proof} environment of 100 lines). No mathematics is written, repaired or re-derived: the statement and the proof are the article's own lines, in the article's own notation, and no retained line is substituted or re-flowed.  Retained uncounted as the source-local context the proof needs: lines 238--242.  Nothing was omitted from any retained span, so the package carries no omission record.  No retained line contains a citation, so the wrapper carries no bibliography and no bibliography entry.  No retained line contains a figure environment, an image-inclusion command or drawing code, so no image is delivered and the package declares no asset dependency.  Editorial formatting only, in addition: an AMS \texttt{amsart} preamble that restates the article's own declarations and macros used by the retained lines, namely the environments \texttt{theorem} on separate counters, together with the standard packages those lines need.  The retained environments print in this document as Theorem 1 in order of appearance, whereas the article prints the same items with the numbering of its own full document.  The complete native source of the article as distributed in the arXiv e-print for this exact version is bundled unchanged as \texttt{sources/189520/source0.tex}.
\begin{equation}\label{eq:qpoch}
(q)_\infty
  \coloneqq
  \prod_{j=1}^{\infty} (1 - q^j).
\end{equation}

\begin{theorem}\label{thm:euler-factorization} We have 
\begin{equation}
    \mathcal{F}(q)
= \frac{1}{(q)_\infty}
\left(
\widehat{F}_0^{(\infty)}(q) + \widehat{F}_1^{(\infty)}(q)
\right), 
\end{equation}
where the modified sequences $\widehat{F}_0^{(n)}$ and $\widehat{F}_1^{(n)}$ satisfy the following coupled recurrences  
\begin{align}
\widehat{F}_0^{(n)} &= \widehat{F}_0^{(n-1)} + q^n\widehat{F}_1^{(n-1)},\label{eq:FirstRec}\\
\widehat{F}_1^{(n)} &= q^n \widehat{F}_0^{(n-1)} + (1-q^n)\widehat{F}_1^{(n-1)},\label{eq:SecondRec}
\end{align}
with initial conditions \(\widehat{F}_0^{(0)}=1,\ \widehat{F}_1^{(0)}=0.\)
\end{theorem}

\begin{proof}

% {\color{red}
% \begin{align*}
% \widehat{F}_0^{(n)}-\widehat{F}_1^{(n-1)} &= \widehat{F}_0^{(n-1)} ,\\
% \widehat{F}_1^{(n)} &= q^n \widehat{F}_0^{(n)} + (1-2q^n)\widehat{F}_1^{(n-1)},
% \end{align*}
% }

 

The upper-left entry of the transfer matrix $M_j(q)$ is simply a   classical Euler factor $\dfrac{1}{1-q^j}$, which is precisely the \(j\)-th factor in the Euler product (or the infinite $q$-Pochhammer symbol) \eqref{eq:qpoch}.


Extracting these factors from the infinite product
\(\prod_{j\ge1} M_j(q)\) leads to a clean normalization of the transfer-matrix
recurrence.



 
To make this precise, define the \emph{normalized} matrices
\begin{equation}\label{eq:NormlizedMatrix}
    \widehat{M}_j(q)
 = 
(1-q^j)\, M_j(q)
=
\begin{pmatrix}
1 & q^j \\[4pt]
q^j & 1-q^j
\end{pmatrix}.
\end{equation}
Then for any row vector \(V=(V_0,V_1)\) we have
\[
V M_j(q)
=
\frac{1}{1-q^j}\, V \widehat{M}_j(q).
\]
It follows that
\[
(1,0)\!\left(\prod_{j\ge1} M_j(q)\right)
=
\left(\prod_{j\ge1} \frac{1}{1-q^j}\right)
(1,0)\!\left(\prod_{j\ge1} \widehat{M}_j(q)\right)
=
\frac{1}{(q)_\infty}\,
(1,0)\!\left(\prod_{j\ge1} \widehat{M}_j(q)\right).
\]



 
Let \(\bigl(\widehat{F}_0^{(n)}(q),\widehat{F}_1^{(n)}(q)\bigr)\) denote the row vector
obtained after multiplying the normalized matrices \(\widehat{M}_1(q),\ldots,\widehat{M}_n(q)\), that is, the \emph{normalized row-vectors} 
\begin{equation}\label{eq:NormlizedVector}
    \bigl(\widehat{F}_0^{(n)},\widehat{F}_1^{(n)}\bigr)
=
(1,0) \,\widehat{M}_1(q)\,\widehat{M}_2(q)\cdots\widehat{M}_n(q).
\end{equation}
Therefore, we have the recurrence
\begin{equation}
    \bigl(\widehat{F}_0^{(n)},\widehat{F}_1^{(n)}\bigr)
=
\bigl(\widehat{F}_0^{(n-1)},\widehat{F}_1^{(n-1)}\bigr)
\begin{pmatrix}
1 & q^n \\[4pt]
q^n & 1-q^n
\end{pmatrix},
\end{equation}
which yields the component-wise system
\begin{align*}
\widehat{F}_0^{(n)}
&=
\widehat{F}_0^{(n-1)} + q^n\,\widehat{F}_1^{(n-1)},\\[4pt]
\widehat{F}_1^{(n)}
&=
q^n\,\widehat{F}_0^{(n-1)} + (1-q^n)\,\widehat{F}_1^{(n-1)}.
\end{align*}
The initial conditions come from \((1,0)\) before any matrices are applied:
\[
\widehat{F}_0^{(0)} = 1,\qquad
\widehat{F}_1^{(0)} = 0.
\]



 
Finally, summing over both terminal states gives
\[
\mathcal{F}(q)
=
\begin{pmatrix}1&0\end{pmatrix}
\!\left(\prod_{j\ge1} M_j(q)\right)
\begin{pmatrix}1\\1\end{pmatrix}
=
\frac{1}{(q)_\infty}\,
\Bigl(\widehat{F}_0^{(\infty)}(q)+\widehat{F}_1^{(\infty)}(q)\Bigr),
\]
proving the normalized Euler-type factorization.
\end{proof}

\end{document}
